\documentclass[11pt,a4paper]{article}
\usepackage{iftex}
\ifPDFTeX
\usepackage[T1]{fontenc}\usepackage{lmodern}
\pdfsuppressptexinfo=15
\fi
\usepackage{amsmath,amssymb}
\usepackage[margin=23mm]{geometry}
\usepackage[hidelinks,bookmarks=false]{hyperref}
\usepackage{microtype}
\setlength{\parindent}{0pt}
\setlength{\parskip}{5pt}
\setlength{\emergencystretch}{2em}
\raggedbottom
\newcommand{\heading}[1]{\par\vspace{5pt}{\large\bfseries #1}\par\vspace{1pt}}
\hypersetup{pdftitle={Bulk Ising susceptibility: a candidate natural-boundary proof},
 pdfauthor={Dehao Lin},
 pdfsubject={Expert brief for manuscript v0.1-rc4 and Lean formalization / post-audit release v0.2.1; AI-assisted candidate proof}}

\begin{document}
{\Large\bfseries Bulk Ising susceptibility: a candidate natural-boundary proof}\par
{\small Expert brief for manuscript v0.1-rc4; Lean formalization / post-audit release v0.2.1}\par
Dehao Lin, School of Physics, Sun Yat-sen University\\
{\small ORCID: \href{https://orcid.org/0009-0001-4551-8490}{0009-0001-4551-8490}}\\
\texttt{lindh9@mail2.sysu.edu.cn}\hfill 7 October 2026\\
{\small Repository: \url{https://github.com/LeoLam233/ising-bulk-natural-boundary}}

\heading{The object and the gap}
The accompanying manuscript presents a candidate complete proof that \(|s|=1\) is a
natural boundary for the zero-field, isotropic, infinite square-lattice
\textbf{bulk susceptibility}, continued from the low-temperature \( + \) boundary-condition state.
This candidate is prepared for public mathematical review.
For \(s=\sinh(2\beta_{\rm phys}J)\), the normalized exterior germ inherited
from real \(s>1\) is
\[
\mathcal X(s)=\beta_{\rm phys}^{-1}\chi(s)
=1-\mathcal M^2(s)+2\mathcal M^2(s)
 \sum_{\substack{N\ge2\\N\ {\rm even}}}T_N(s),
\qquad \mathcal M^2(s)=(1-s^{-4})^{1/4}.
\]
The response sums connected correlations over the full lattice and retains
the infinite particle expansion.

Nickel (1999, 2000) identified fixed-form-factor singularities whose union
is dense on the unit circle. Orrick et al. (2001) gave amplitude-based
evidence against cancellation in the full sum, explicitly short of a
proof. Tracy--Widom (2014) proved smooth extension
away from Nickel points at each fixed even order. Their diagonal-susceptibility
(2013) and Toeplitz-sum (2015/2017) works provide rigorous noncancellation
precedents. The gap addressed here is control of the
\emph{entire higher-particle tail in the bulk response} near the boundary,
so that it cannot cancel the first singular contribution.

\heading{Two estimates that carry the conclusion}
For each fixed point \(s_*\) in the selected cyclotomic family of
\S3, let \(N_0=2p\) be its first even order, with \(p\ge11\) prime,
\(k=N_0^2/2-1\), and \(s_\varepsilon=(1+\varepsilon)s_*\).
The two claimed estimates, as \(\varepsilon\downarrow0\), are
\begin{align*}
\partial_s^kT_{N_0}(s_\varepsilon)
 &=2L_\beta\varepsilon^{-1/2}+o(\varepsilon^{-1/2}),
 \qquad L_\beta\ne0,\\[5pt]
\sum_{\substack{N>N_0\\N\ {\rm even}}}
 \bigl|\partial_s^jT_N(s_\varepsilon)\bigr|
 &=o(\varepsilon^{-1/2}),\qquad 0\le j\le k.
\end{align*}
\textbf{The first supplies a nonzero singular carrier; the second makes
the entire higher-particle tail too small to cancel it.}
The first coefficient belongs to the whole form factor, not just a local
model. The tail covers every higher even order and the full angular domain.
Constants may depend on the fixed point and finite \(k\). With the stated
lower-order bounds, normal convergence and symmetries, these estimates
yield the claimed full natural boundary (\S9).

\heading{The proposed bulk implementation}
The technical claims are a dense uniquely-first point family, a
whole-form-factor first-amplitude calculation, an exact weighted
contour/Stokes decomposition, protected parameter disks, and differentiated
full-tail/coarea estimates. The noncancellation principle is established
prior work. Contour partitions, Vandermonde analysis, grouping and
Hadamard/factorial budgets also have Tracy--Widom precedents (\S1.1).
No exhaustive priority clearance is claimed.

\textbf{External inputs.} Palmer--Tracy's (1981) low-temperature Fredholm
correlation representation, as used in the published Tracy--Widom (2014), Appendix 1;
Yang's (1952) magnetization; and Tracy--Widom's fixed-even-order non-Nickel smoothness.
None supplies the bulk infinite-tail theorem. Sources and conventions
are in manuscript \S1.1, \S2 and Appendix F.

\newpage
\heading{Proof architecture}
The two branches of the argument meet in the noncancellation step:
\[
\begin{aligned}
\text{Dense uniquely-first family}
 &\longrightarrow \text{nonzero whole first singularity},\\[3pt]
T_N=F_N+K_N+S_N
 &\longrightarrow \text{uniform differentiated full tail},\\[3pt]
\text{Both estimates}
 &\longrightarrow \text{dense boundary noncancellation}\\[-1pt]
 &\longrightarrow \text{natural boundary}.
\end{aligned}
\]
Here \(F_N\) is the final selected-contour integral, \(K_N\) the retained
original contribution, and \(S_N\) the weighted Stokes correction.
The tail argument combines protected-disk estimates, pair/Pfaffian
suppression and differentiation/coarea bounds near collisions.

\heading{Suggested independent reading}
Each interface below can be read separately as a first check.

\textbf{1. First carrier and amplitude --- \S4, Appendix B.}
Does the fixed \(y\)-only localization, followed by full \(x\) residues
and the mean-angle deformation, give the coefficient of the \emph{whole}
form factor? Check that the two chart coefficients add without cancellation
and that the complement and lower derivatives remain controlled.

\textbf{2. Weighted contours and pair control --- \S\S5--8.}
Check the exact identity \(T_N=F_N+K_N+S_N\) on the actual coupled
support, including its Stokes correction. Do the protected parameter disks,
regularized complete pairs and same-group contraction survive branch and
endpoint degeneracies? Does the selected-contour Pfaffian budget retain
every differentiation and integration cost?

\textbf{3. High-order differentiation and coarea --- Appendices D--E.}
Does the Lie-derivative recurrence include derivatives of previously
generated coefficients and real cutoffs? Check puncture flux at each stage,
the real coarea Jacobian and multiplicity, and the order of constants and
quantifiers needed for the uniform estimates.

\textbf{4. Complete tail closure --- \S9, Appendix G.}
Do the thirteen sectors W1--W13 cover all higher even orders and all
\(0\le j\le k\), including transitions between supports and particle-number
windows? Verify that their bounds sum to the stated smaller-order tail.

\heading{Specific feedback sought}
A concrete failing configuration, an unjustified implication or uniformity
step, or a missed close precedent would be especially useful.
\textbf{A limited reading of one interface is valuable; a full refereeing
commitment is not presumed.}

\textbf{Current verification status.} The internal form-factor and higher-tail proof chain is
formalized in Lean 4. An independent adversarial audit of frozen v0.2.0 gave
\texttt{PASS\_WITH\_NONBLOCKING\_ISSUES}, finding no blocking or load-bearing defect.
Release v0.2.1 applied targeted post-audit remediation without changing the core
natural-boundary theorems. The exact v0.2.1 commit passed the final targeted
regressions and controls, a clean full build, a project-wide axiom audit of 9,200
safe declarations, 14 endpoint checks, fresh kernel replay, and full and smoke CI.
E1/E2/E3 remain explicit external literature inputs: the low-temperature correlation
representation, Yang's spontaneous-magnetization formula, and Tracy--Widom's
fixed-even-order non-Nickel smoothness. This remains an AI-assisted candidate
proof: no independent human expert validation or peer review is claimed.

\vfill
{\small Section and appendix pointers refer to the accompanying
manuscript v0.1-rc4; full bibliographic details are in its reference list.\\
Project text: CC BY 4.0. Project code: MIT. Cited works retain their own rights.}
\end{document}
